% =====================================================================
% DRAFT -- standalone derivation sketch. NOT main.tex.
% Purpose: recover the PoLoRA direction from a per-factor layerwise-KL model.
% =====================================================================
\documentclass[11pt]{article}
\usepackage{amsmath,amssymb,amsthm,mathtools}
\usepackage[margin=1in]{geometry}

\DeclarePairedDelimiter{\norm}{\lVert}{\rVert}
\newcommand{\polar}{\operatorname{polar}}
\newcommand{\DW}{\Delta W}
\newcommand{\R}{\mathbb{R}}
\newcommand{\E}{\mathbb{E}}
\newcommand{\T}{\mathcal T}
\newcommand{\F}{\mathcal F}
\DeclareMathOperator{\tr}{tr}
\DeclareMathOperator*{\argmin}{arg\,min}
\DeclareMathOperator{\diag}{diag}

\newtheorem{proposition}{Proposition}
\newtheorem{lemma}{Lemma}
\theoremstyle{remark}
\newtheorem*{remark}{Remark}

\title{PoLoRA from a per-factor layerwise-KL model\\[2pt]
\large (derivation draft --- for discussion, not for \texttt{main.tex})}
\author{}
\date{}

\begin{document}
\maketitle

This note derives the PoLoRA update directly in LoRA factor coordinates
$(A,B)$. The model is a \emph{layerwise-KL proximal} problem: a first-order
task-loss term plus a second-order penalty that measures how far the step moves
the layer's predictive distribution. Solving it separately for each factor
produces PoLoRA's whiten--polar--unwhiten step \eqref{eq:final-step}.

The derivation runs in two stages.
\begin{itemize}
  \item A dense one-variable \emph{warmup} (\S\ref{sec:dense-warmup}) on a
        hypothetical dense update $U\in\R^{d_{\mathrm{out}}\times d_{\mathrm{in}}}$
        to the full adapter matrix $W=BA$. It carries one update variable and one
        damping, so it is \emph{not} the PoLoRA step; its only job is to show why
        a layerwise-KL metric plus a worst-input penalty yields a
        whiten--polar--unwhiten direction.
  \item The \emph{per-factor model} (\S\ref{sec:channel-model} onward), where the
        two parts of the tangent update, $B\Delta A$ and $\Delta B\,A$, each carry
        their own damping $\lambda_A,\lambda_B$ from the start. This is the actual
        derivation.
\end{itemize}
\noindent\S\ref{sec:local-kl} sets up the layerwise-KL metric;
\S\ref{sec:worst-input}--\S\ref{sec:radius} solve the two per-factor subproblems
and calibrate the dampings; \S\ref{sec:est} estimates the metric.

\section{Setup}

A LoRA layer adds the rank-$r$ adapter
$W=BA\in\R^{d_{\mathrm{out}}\times d_{\mathrm{in}}}$ to a frozen weight, with
$A\in\R^{r\times d_{\mathrm{in}}}$ and
$B\in\R^{d_{\mathrm{out}}\times r}$. One optimizer step increments the factors by
$\Delta A\in\R^{r\times d_{\mathrm{in}}}$ and
$\Delta B\in\R^{d_{\mathrm{out}}\times r}$, changing $W$ by
\begin{equation}
  U:=B\,\Delta A+\Delta B\,A+\Delta B\,\Delta A
   \approx B\,\Delta A+\Delta B\,A .
  \label{eq:tangent}
\end{equation}
Let $G=\nabla_W\mathcal L\in\R^{d_{\mathrm{out}}\times d_{\mathrm{in}}}$ be the
gradient with respect to $W$. The first-order loss change splits across the two
factors,
\begin{equation}
  \langle G,U\rangle
  =\langle B^\top G,\Delta A\rangle
  +\langle G A^\top,\Delta B\rangle ,
  \qquad
  G_A:=B^\top G\in\R^{r\times d_{\mathrm{in}}},\quad
  G_B:=G A^\top\in\R^{d_{\mathrm{out}}\times r} .
  \label{eq:split-linear}
\end{equation}
The derivation uses raw gradients; the optimizer replaces $G_A,G_B$ by
bias-corrected momentum buffers before the polar map.

\section{Layerwise KL supplies the layer metric}
\label{sec:local-kl}

Newton--Muon derives its triplet quadratic by Taylor expanding the task loss in a
dense layer update. PoLoRA uses a different model: the task loss is kept to first
order, while the second-order term is the predictive KL. The algebra has the same
three ingredients---a gradient, layer inputs, and an output-side metric---but the
metric is layerwise-KL/Fisher geometry rather than the observed-loss Hessian.

Let $z_i\in\R^{d_{\mathrm{in}}}$ be the layer input for token or example $i$,
collect $Z=[z_1,\ldots,z_N]$, and let
$y_i=Wz_i\in\R^{d_{\mathrm{out}}}$ be the layer output. For token $i$, let
$p_i(\cdot;y)$ be the model's predictive distribution when the layer output is
$y$. For an output perturbation $v_i$, define
\[
  D_i(v_i)
  :=
  \mathrm{KL}\!\left(p_i(\cdot;y_i)\,\middle\|\,p_i(\cdot;y_i+v_i)\right).
\]

\begin{lemma}[Layerwise-KL expansion]\label{lem:local-kl}
Assume the usual regularity conditions allowing differentiation under the
expectation. Let
\[
  s_i(a):=\nabla_y \log p_i(a;y_i),
  \qquad
  F_i:=\E_{a\sim p_i(\cdot;y_i)}[s_i(a)s_i(a)^\top].
\]
Then
\[
  D_i(v_i)=\tfrac12 v_i^\top F_i v_i+o(\norm{v_i}_2^2).
\]
\end{lemma}

\begin{proof}
The first derivative at $v_i=0$ is
$-\E_{a\sim p_i(\cdot;y_i)}[\nabla_y\log p_i(a;y_i)]$, which is zero by the score
identity. The second derivative is
$-\E_{a\sim p_i(\cdot;y_i)}[\nabla_y^2\log p_i(a;y_i)]$, which equals $F_i$ by
the Fisher identity. Taylor expansion gives the claim.
\end{proof}

\begin{lemma}[When Fisher equals Hessian]\label{lem:fisher-hessian}
If the token loss is negative log-likelihood, $\ell_i(y)=-\log p_i(a;y)$, and the
label $a$ is averaged under the current model distribution $p_i(\cdot;y_i)$, then
\[
  \E_{a\sim p_i(\cdot;y_i)}[\nabla_y^2\ell_i(y_i)]
  =
  F_i .
\]
\end{lemma}

\begin{proof}
Since $\ell_i(y)=-\log p_i(a;y)$,
\[
  \E[\nabla_y^2\ell_i(y_i)]
  =
  -\E[\nabla_y^2\log p_i(a;y_i)]
  =
  \E[s_i(a)s_i(a)^\top],
\]
using the same Fisher identity as Lemma~\ref{lem:local-kl}.
\end{proof}

Thus Fisher is an exact expected Hessian only under the model-distribution
expectation in Lemma~\ref{lem:fisher-hessian}. For observed labels in a finite
minibatch, the empirical gradient outer product is a PSD estimator or proxy for
the layerwise-KL metric, not an identity with the raw loss Hessian.

\section{Dense warmup: one update, one damping}
\label{sec:dense-warmup}

The warmup carries one dense update variable $U$ and one damping $\lambda>0$, so
it is not the PoLoRA factor step. Its only purpose is to show why a layerwise-KL
metric combined with a worst-input replacement produces a polar direction.

For a hypothetical dense update $U\in\R^{d_{\mathrm{out}}\times d_{\mathrm{in}}}$
to $W$, the output perturbation on input $z_i$ is $Uz_i$. The first-order
task-loss term plus the second-order layerwise-KL expansion gives
\begin{equation}
  \mathcal J_{\mathrm{dense}}(U)
  =
  \langle G,U\rangle
  +\frac{\lambda}{2N}\sum_{i=1}^N (Uz_i)^\top F_i(Uz_i).
  \label{eq:dense-local-kl}
\end{equation}
The corresponding unscaled dense layerwise-KL metric is
\begin{equation}
  \Sigma_{\mathrm{KL}}
  =
  \frac1N\sum_{i=1}^N (z_i z_i^\top)\otimes F_i .
  \label{eq:param-kl}
\end{equation}
PoLoRA never forms this dense matrix. The KL-Shampoo step instead fits
$\Sigma_{\mathrm{KL}}$ (or its stochastic gradient-covariance estimator) by a
diagonal--Kronecker matrix
\begin{equation}
  \Sigma_{\mathrm{KL}}\approx Q\otimes P,
  \label{eq:PQ-approx}
\end{equation}
where $P\in\R^{d_{\mathrm{out}}\times d_{\mathrm{out}}}$ is the output-side factor
and $Q\in\R^{d_{\mathrm{in}}\times d_{\mathrm{in}}}$ is the input-side factor, both
symmetric positive definite, with their coupled estimator given in
\S\ref{sec:est}. So $P,Q$ are the diagonal Kronecker fit to the layer's
layerwise-KL metric, not dense metric factors. The scale ambiguity
$(P,Q)\mapsto(cP,c^{-1}Q)$ is harmless because $P^{1/2}UQ^{1/2}$ is unchanged.

Replacing \eqref{eq:param-kl} by \eqref{eq:PQ-approx} gives the fitted average
model
\begin{equation}
  \mathcal J_{\mathrm{avg}}(U)
  =
  \langle G,U\rangle
  +\frac{\lambda}{2}\norm{P^{1/2}UQ^{1/2}}_F^2 .
  \label{eq:pq-avg}
\end{equation}
Replacing the average squared response by the largest one over unit inputs gives
the worst-input model
\begin{equation}
  \mathcal J_{\mathrm{worst}}(U)
  =
  \langle G,U\rangle
  +\frac{\lambda}{2}\norm{P^{1/2}UQ^{1/2}}_2^2
  =
  \langle G,U\rangle
  +\frac{\lambda}{2}
    \sup_{\norm{x}_2\le1}\norm{P^{1/2}UQ^{1/2}x}_2^2 .
  \label{eq:robust-pq}
\end{equation}
Let $Y=P^{1/2}UQ^{1/2}$ and $H=P^{-1/2}GQ^{-1/2}$. Then
\begin{equation}
  \min_U \mathcal J_{\mathrm{worst}}(U)
  =
  \min_Y
  \left\{
    \langle H,Y\rangle
    +\frac{\lambda}{2}\norm{Y}_2^2
  \right\}.
  \label{eq:dense-penalty}
\end{equation}
By the scaled conjugacy between squared operator norm and squared nuclear norm,
a canonical minimizer, when $H\ne0$, is
\[
  Y_\star=-\frac{\norm{H}_*}{\lambda}\polar(H).
\]
Thus the dense warmup gives
\begin{equation}
  U_\star
  =
  -\frac{\norm{H}_*}{\lambda}\,
  P^{-1/2}\polar\!\big(P^{-1/2}GQ^{-1/2}\big)Q^{-1/2}.
  \label{eq:merged-polar}
\end{equation}
The lesson is only the direction mechanism: layerwise-KL whitening plus the
worst-input replacement gives a whiten--polar--unwhiten update. The scalar
$\norm{H}_*/\lambda$ belongs to the dense one-variable proximal problem and is
not the PoLoRA factor magnitude. PoLoRA starts from the per-factor model below.

\section{Per-factor proximal model}
\label{sec:channel-model}

A LoRA tangent update splits into two parts,
\[
  J_A:=B\Delta A,\qquad
  J_B:=\Delta B A,\qquad
  J=J_A+J_B,
\]
the part driven by $\Delta A$ and the part driven by $\Delta B$. For input $z_i$,
the two output perturbations are
\[
  v_{A,i}:=B\Delta A z_i,\qquad
  v_{B,i}:=\Delta B A z_i.
\]
The exact second-order layerwise KL of the simultaneous tangent perturbation
would contain
\[
  \frac1{2N}\sum_{i=1}^N
  (v_{A,i}+v_{B,i})^\top F_i (v_{A,i}+v_{B,i}),
\]
including the cross terms $v_{A,i}^\top F_i v_{B,i}$. PoLoRA makes a
block-diagonal approximation across the two parts and assigns one proximal
damping to each:
\begin{equation}
  \Phi_{\lambda_A,\lambda_B}(\Delta A,\Delta B)
  =
  \langle G_A,\Delta A\rangle+\langle G_B,\Delta B\rangle
  +\frac{\lambda_A}{2N}\sum_{i=1}^N v_{A,i}^\top F_i v_{A,i}
  +\frac{\lambda_B}{2N}\sum_{i=1}^N v_{B,i}^\top F_i v_{B,i}.
  \label{eq:channel-prox}
\end{equation}
This is still a proximal layerwise-KL model; no penalty has been converted into a
constraint. The two multipliers are present from the start because the retained
metric is block diagonal across the $\Delta A$ and $\Delta B$ blocks. If the
cross-term metric block were kept and the two dampings were tied, the model would
reduce to the exact layerwise-KL quadratic for $J_A+J_B$ under a common damping.

Using the KL-Shampoo fit \eqref{eq:PQ-approx}, the $\Delta A$ part's layerwise-KL
quadratic becomes
\begin{equation}
  \frac1N\sum_{i=1}^N (B\Delta A z_i)^\top F_i(B\Delta A z_i)
  \approx
  \norm{P^{1/2}B\Delta A Q^{1/2}}_F^2
  =
  \norm{C_A^{1/2}\Delta A Q^{1/2}}_F^2,
  \label{eq:A-average-fit}
\end{equation}
where
\begin{equation}
  C_A:=B^\top P B\in\R^{r\times r} .
  \label{eq:CA}
\end{equation}
The $\Delta B$ part's quadratic becomes
\begin{equation}
  \frac1N\sum_{i=1}^N (\Delta B A z_i)^\top F_i(\Delta B A z_i)
  \approx
  \norm{P^{1/2}\Delta B A Q^{1/2}}_F^2
  =
  \norm{P^{1/2}\Delta B C_B^{1/2}}_F^2,
  \label{eq:B-average-fit}
\end{equation}
where
\begin{equation}
  C_B:=AQA^\top\in\R^{r\times r} .
  \label{eq:CB}
\end{equation}
The equalities in \eqref{eq:A-average-fit} and \eqref{eq:B-average-fit} are
ordinary trace identities. For example,
\[
  \tr(PB\Delta A Q\Delta A^\top B^\top)
  =
  \tr(C_A\Delta A Q\Delta A^\top).
\]

\section{Worst-input proximal problems}
\label{sec:worst-input}

The KL-Shampoo fit gives an average squared response. PoLoRA replaces that
average by the largest response over unit inputs, separately for each factor:
\begin{equation}
  R_A(\Delta A)
  :=
  \norm{C_A^{1/2}\Delta A Q^{1/2}}_2^2
  =
  \sup_{\norm{x}_2\le1}
  \norm{C_A^{1/2}\Delta A Q^{1/2}x}_2^2,
  \label{eq:RA}
\end{equation}
\begin{equation}
  R_B(\Delta B)
  :=
  \norm{P^{1/2}\Delta B C_B^{1/2}}_2^2 .
  \label{eq:RB}
\end{equation}
Substituting \eqref{eq:RA}--\eqref{eq:RB} into the per-factor model gives two
separable subproblems:
\begin{equation}
  \mathcal J_A(\Delta A)
  =
  \langle G_A,\Delta A\rangle
  +\frac{\lambda_A}{2}
    \norm{C_A^{1/2}\Delta A Q^{1/2}}_2^2,
  \label{eq:A-block}
\end{equation}
\begin{equation}
  \mathcal J_B(\Delta B)
  =
  \langle G_B,\Delta B\rangle
  +\frac{\lambda_B}{2}
    \norm{P^{1/2}\Delta B C_B^{1/2}}_2^2 .
  \label{eq:B-block}
\end{equation}

\begin{proposition}[Per-factor proximal solutions]\label{prop:channels}
Assume the damped matrices $P,Q,C_A,C_B$ are positive definite, and set
\[
  H_A:=C_A^{-1/2}G_AQ^{-1/2},
  \qquad
  H_B:=P^{-1/2}G_BC_B^{-1/2}.
\]
If $H_A\ne0$ and $H_B\ne0$, canonical minimizers of
\eqref{eq:A-block}--\eqref{eq:B-block} are
\begin{equation}
  \Delta A(\lambda_A)
  =
  -\frac{\norm{H_A}_*}{\lambda_A}\,
  C_A^{-1/2}\polar(H_A)Q^{-1/2},
  \label{eq:A-solution}
\end{equation}
\begin{equation}
  \Delta B(\lambda_B)
  =
  -\frac{\norm{H_B}_*}{\lambda_B}\,
  P^{-1/2}\polar(H_B)C_B^{-1/2}.
  \label{eq:B-solution}
\end{equation}
If a whitened gradient is zero, the corresponding zero update is a minimizer.
\end{proposition}

\begin{proof}
With $Y_A=C_A^{1/2}\Delta A Q^{1/2}$,
$\langle G_A,\Delta A\rangle=\langle H_A,Y_A\rangle$
and the quadratic term is $\frac{\lambda_A}{2}\norm{Y_A}_2^2$. By the scaled
conjugacy between squared operator norm and squared nuclear norm, the minimum
value is $-\norm{H_A}_*^2/(2\lambda_A)$. Since
$\polar(H_A)\in\partial\norm{H_A}_*$, the canonical
attainer is
$Y_A=-(\norm{H_A}_*/\lambda_A)\polar(H_A)$. Undoing the
change of variables gives \eqref{eq:A-solution}. The $\Delta B$ proof is
identical with $Y_B=P^{1/2}\Delta B C_B^{1/2}$.
\end{proof}

Define the unscaled update directions
\begin{equation}
  W_A:=C_A^{-1/2}\polar(H_A)Q^{-1/2},
  \label{eq:WA}
\end{equation}
\begin{equation}
  W_B:=P^{-1/2}\polar(H_B)C_B^{-1/2}.
  \label{eq:WB}
\end{equation}

\section{Damping calibration from the product radius}
\label{sec:radius}

The proximal solutions \eqref{eq:A-solution}--\eqref{eq:B-solution} have
operator norms
\begin{equation}
  \norm{\Delta A(\lambda_A)}_2
  =
  \frac{\norm{H_A}_*}{\lambda_A}\norm{W_A}_2,
  \qquad
  \norm{\Delta B(\lambda_B)}_2
  =
  \frac{\norm{H_B}_*}{\lambda_B}\norm{W_B}_2 .
  \label{eq:block-norms}
\end{equation}
Thus $\lambda_A$ and $\lambda_B$ control only the scalar lengths along the
layerwise-KL polar directions. PoLoRA chooses these dampings so the factor
updates have a common operator-norm radius
\begin{equation}
  \norm{\Delta A}_2=\norm{\Delta B}_2=\rho .
  \label{eq:equal-radius}
\end{equation}
The first-order tangent update then obeys
\begin{equation}
  \norm{B\Delta A+\Delta B A}_2
  \le
  \norm{B}_2\norm{\Delta A}_2+\norm{A}_2\norm{\Delta B}_2
  =\rho(\norm{A}_2+\norm{B}_2).
\end{equation}
Setting this bound equal to the budget $\eta$ gives
\begin{equation}
  \rho=\frac{\eta}{\norm{A}_2+\norm{B}_2},
  \label{eq:rho}
\end{equation}
and the dampings that realize this radius are
\begin{equation}
  \lambda_A
  =
  \frac{\norm{H_A}_*\norm{W_A}_2}{\rho},
  \qquad
  \lambda_B
  =
  \frac{\norm{H_B}_*\norm{W_B}_2}{\rho}.
  \label{eq:lambdas}
\end{equation}
Substituting \eqref{eq:lambdas} into
\eqref{eq:A-solution}--\eqref{eq:B-solution} gives the PoLoRA step
\begin{equation}
  \Delta A=-\rho\,\frac{W_A}{\norm{W_A}_2},
  \qquad
  \Delta B=-\rho\,\frac{W_B}{\norm{W_B}_2}.
  \label{eq:final-step}
\end{equation}
for nonzero $W_A,W_B$.

\section{Estimating the diagonal metric}
\label{sec:est}

The derivation assumes $P,Q$ fixed during the step. PoLoRA estimates them by the
coupled diagonal KL-Shampoo moment equations. This estimator belongs to the same
layerwise-KL modelling choice, but it does not require forming the dense metric
\eqref{eq:param-kl}. Instead it uses the stochastic dense layer gradient as an
empirical-Fisher estimator of that metric:
\begin{equation}
  \widehat\Sigma_{\mathrm{KL}}
  :=
  \E[\mathrm{vec}(G)\mathrm{vec}(G)^\top]
  \approx Q\otimes P .
  \label{eq:efisher-kron}
\end{equation}
For likelihood losses under the model-distribution expectation,
Lemma~\ref{lem:fisher-hessian} identifies this gradient second moment with the
layerwise-KL metric. With observed labels, it is the empirical-Fisher proxy for
the same metric. In both cases, KL-Shampoo replaces the dense metric by the
diagonal--Kronecker projection that best matches the observed gradient
covariance.

The only algebraic consequences needed from \eqref{eq:efisher-kron} are the moment
identities, for compatible $M,N$,
\begin{equation}
  \E[G M G^\top]=\tr(MQ)P,
  \qquad
  \E[G^\top N G]=\tr(NP)Q .
  \label{eq:moment-identities}
\end{equation}
The factor gradients are linear images of $G$, so their second moments carry the
same $r\times r$ matrices $C_A,C_B$ that appear in the update:
\begin{equation}
  \E[\mathrm{vec}(G_A)\mathrm{vec}(G_A)^\top]\approx Q\otimes C_A,
  \qquad
  \E[\mathrm{vec}(G_B)\mathrm{vec}(G_B)^\top]\approx C_B\otimes P .
\end{equation}

\begin{lemma}[Coupled diagonal moment estimator]\label{lem:estimator}
With $C_A=B^\top P B$ and $C_B=AQA^\top$,
\begin{equation}
  \E\!\left[\frac1r G_A^\top C_A^{-1}G_A\right]=Q,
  \qquad
  \E\!\left[\frac1r G_B C_B^{-1}G_B^\top\right]=P .
  \label{eq:unbiased}
\end{equation}
\end{lemma}

\begin{proof}
Use the second identity of \eqref{eq:moment-identities} with $N=B C_A^{-1}B^\top$:
$G^\top N G=G_A^\top C_A^{-1}G_A$ and
$\tr(NP)=\tr(C_A^{-1}B^\top P B)=\tr(I_r)=r$, giving the first identity. The second
follows by the same argument with $M=A^\top C_B^{-1}A$.
\end{proof}

PoLoRA keeps only the diagonals and updates them by EMA:
\begin{equation}
  q_j\leftarrow
  \mathrm{EMA}\!\left[\frac1r (G_A)_{:,j}^\top C_A^{-1}(G_A)_{:,j}\right],
  \qquad
  p_i\leftarrow
  \mathrm{EMA}\!\left[\frac1r (G_B)_{i,:}C_B^{-1}(G_B)_{i,:}^\top\right].
  \label{eq:ema}
\end{equation}
Because $C_A$ depends on the current $P$ and $C_B$ on the current $Q$, the two
estimates are coupled; the optimizer takes one EMA alternation per step, then
freezes $P,Q$ while computing \eqref{eq:WA}--\eqref{eq:final-step}. In
implementation the inverse square roots are damped, so the formulas are read as the
ideal SPD case.

\paragraph{Summary.}
Starting from the per-factor layerwise-KL proximal model \eqref{eq:channel-prox},
KL-Shampoo fits the retained metric by $Q\otimes P$, the worst-input replacement
turns each penalty into a spectral norm \eqref{eq:RA}--\eqref{eq:RB}, and the
per-factor proximal solutions are the polar sandwiches
\eqref{eq:WA}--\eqref{eq:WB}. Calibrating $\lambda_A,\lambda_B$ to the common
product radius \eqref{eq:rho} fixes the factor lengths and gives the implemented
step \eqref{eq:final-step}. The dense warmup
\eqref{eq:dense-local-kl}--\eqref{eq:merged-polar} only motivates the polar
mechanism; its single damping never enters the per-factor model.

\end{document}
